\chapter{梯度提升树及其应用}
\chaptergoal{理解依次加入弱学习器改进整体预测的思想；从平方误差下的残差拟合过渡到一般损失的负梯度；说明学习率、迭代次数和树复杂度的相互关系；使用 sklearn 比较梯度提升与随机森林；了解 XGBoost 的延伸位置。}

随机森林训练许多具有差异的树，再综合它们的预测。梯度提升树（gradient boosting trees）的组织方式不同：先得到一个初始预测，再逐轮增加新树，使整体训练目标朝较低的方向改变。这种顺序依赖意味着某一轮的树要考虑此前模型还没有解释好的部分。它也是本书从第3章损失与梯度概念回到更复杂模型训练的一个重要例子。
\section{逐步改进预测的集成思想}

\textbf{提升（boosting）}是一类按顺序组合模型的策略：后来的模型依据前面模型尚未解决的误差或相应训练目标继续学习。这里每轮加入的简单模型也常被称为\textbf{弱学习器}，本章通常用浅回归树实现。这个称呼强调每一轮修正的表达能力有限，不表示任何随意选出的坏模型组合起来都会成功。

模型总预测是一串修正的\textbf{相加}，而不是随机森林那样把多棵完整预测简单平均。初始预测、每棵修正树及其学习率必须一起保留，预测新样本时仍按同一顺序相加。

考虑回归任务。最简单的初始模型可以为所有样本预测训练目标均值 $F_0(\mathbf x)$。随后训练一棵浅树 $h_1(\mathbf x)$ 作为修正，得到
\begin{equation}
F_1(\mathbf x)=F_0(\mathbf x)+\eta h_1(\mathbf x),
\end{equation}
其中 $\eta$ 是学习率，控制这一轮修正对整体预测的贡献。不断加入新树，就形成 $F_M$。与随机森林的树彼此较独立不同，提升中的新树依赖已有模型的预测情况。

逐步添加的学习器通常相对简单，例如浅回归树。复杂度来自多轮修正的叠加，而不一定要求每一棵树都特别深。这有助于理解为何学习率、迭代轮数和叶节点复杂度需要共同选择。

\subsection{加法模型中的每棵树承担什么角色}
写成一般形式，经过 $M$ 轮的模型为
\[
F_M(\mathbf x)=F_0(\mathbf x)+\sum_{m=1}^M\eta_m h_m(\mathbf x).
\]
$F_0$ 是初始函数，$h_m$ 是第 $m$ 轮新增函数，$\eta_m$控制其贡献。这里的 $h_m$ 通常预测的是修正量，不是对原目标的一份独立完整预测。预测新输入时需要先计算 $F_0$，再累加各轮修正；将它们简单平均，会改变模型的数学含义。

初始函数也由损失确定。平方误差下，常数最优解为训练均值；绝对误差下，可从中位数出发；二分类交叉熵下，常数概率通常对应训练正类比例，其得分是该比例的对数优势。它们与前面章节的常数最优解直接相连，不是所有任务都固定初值为0。

假设训练目标是温度，单棵修正树可以输出负数，因为它表达“原来的预测需要降低多少”，而不是在声称真实温度一定为负。分类提升中，树甚至可能在对数优势得分空间相加，叶值并非类别概率。解释软件输出前，必须先明确相加的量是什么。

\subsection{弱模型组合不是把错误答案重复累加}
如果新树与应修正方向无关，加入它未必降低目标。有效提升需要新学习器能够利用输入信息解释当前梯度结构，同时控制步长和复杂度。模型之间的依赖使后续树围绕前面残留的问题学习，这与Bagging通过重复抽样形成差异的机制不同。

当每轮使用浅树时，单轮只能表达有限的条件交互。多轮相加可形成较丰富函数，但学习率、树深、叶支持和总轮数共同限制表达与泛化。不能只以“树数量很多”来判断森林和提升有相同的容量或训练成本。

\section{平方损失下的残差拟合}

\subsection*{四条记录上的两轮提升}
用人工数据 $x=(1,2,3,4)$、$y=(1,1,3,3)$ 逐步计算，采用平均半平方损失 $J=\frac1{2n}\sum_i(y_i-F_i)^2$。初始预测取目标均值 $F_0=2$，所以四条记录预测为 $(2,2,2,2)$，目标减去预测得到 $(-1,-1,1,1)$。此时 $J_0=0.5$。

第一棵深度为1的回归树可以在 $x\leq2.5$处分裂，左右分别预测修正值$-1$和$1$。取学习率$\eta=0.5$，加上一半修正后，
\[
F_1=(2,2,2,2)+0.5(-1,-1,1,1)=(1.5,1.5,2.5,2.5).
\]
新残差为 $(-0.5,-0.5,0.5,0.5)$，损失降为 $J_1=0.125$。第二棵树拟合这些新残差，仍取一半修正，得到 $F_2=(1.25,1.25,2.75,2.75)$，损失为 $J_2=0.03125$。注意第二轮不是再拟合第一次的残差，而是重新计算当前模型遗漏的部分。

\fig{v2_boost_steps.pdf}{人工四点例子中两轮提升的预测。每轮修正只加入一半；数值可直接与正文计算核对。}

本章用$r=y-F$表示待补偿的残差，与第3章用来写损失导数的$e=\hat y-y$符号相反。两种写法都可用，但这里新树拟合的是$r=-e$，不能直接把方向写反。

若目标是平方误差 $L=\frac12\sum_i(y_i-F(\mathbf x_i))^2$，当前残差 $r_i=y_i-F(\mathbf x_i)$ 指示仍未拟合的部分。下一轮训练一棵树去预测残差，便可能减少整体平方误差。对于初学者，这个例子非常直观：先画出初始预测，再画残差，最后加上一部分残差拟合结果，观察新曲线怎样变化。

\fig{ch10_boosting_steps.pdf}{人工生成的一维回归数据中，初始预测、残差修正和逐轮提升预测的示意。图中的数据是模拟数据，不是科学观测。}

“拟合残差”在这里适用于平方误差的例子，不能作为所有提升算法的唯一解释。若任务是概率分类，训练损失换成适合类别概率的目标，下一轮树所拟合的量会与该目标的梯度有关。

\subsection{当一棵树不能完全解释残差时}
前面的四点例子恰好让一个树桩完整表达当前残差，因此各轮沿同一结构收缩。这是说明学习率的最小示例。更常见的情况是浅树只能近似当前残差，下一轮最优划分也可能改变。

\worked{残差结构改变导致第二棵树改用另一阈值}
取人工输入 $x=(1,2,3,4)$，目标 $y=(0,1,2,4)$。初始均值为1.75，MSE为2.1875。第一轮残差为 $(-1.75,-0.75,0.25,2.25)$。深度1树选择阈值3.5，左侧三个残差均值为 $-0.75$，右侧为2.25。学习率取0.5，得到
\[
F_1=(1.375,1.375,1.375,2.875),\quad \mathrm{MSE}_1=0.921875.
\]
重新计算残差为 $(-1.375,-0.375,0.625,1.125)$。第二个树桩选择阈值2.5，左右均值分别为 $-0.875,0.875$，于是
\[
F_2=(0.9375,0.9375,1.8125,3.3125),\quad\mathrm{MSE}_2=0.34765625.
\]
两个修正树的阈值不同，第一轮只解决了“最后一个目标明显偏高”的主要结构，第二轮再处理前后两部分仍然存在的差异。这里报告的是MSE，不是前一个例子的半MSE；训练方向相同，但数值尺度必须分清。

以下核心循环可复现两个步骤。\texttt{fit}的目标是当前残差，\texttt{predict}后乘学习率才加入整体预测。
\begin{lstlisting}[caption={逐轮重新计算残差，而不是重复使用第一轮残差}]
import numpy as np
from sklearn.tree import DecisionTreeRegressor
X = np.arange(1,5,dtype=float).reshape(-1,1)
y = np.array([0.,1.,2.,4.])
F = np.full(len(y), y.mean())
for stage in range(2):
    residual = y-F
    tree = DecisionTreeRegressor(max_depth=1, random_state=42)
    tree.fit(X, residual)
    F = F + .5*tree.predict(X)
    print(stage+1, F, np.mean((y-F)**2))
\end{lstlisting}

这个循环专门针对平方误差，不应直接替换目标为0/1就宣布实现了概率梯度提升。概率损失还需要说明得分到概率的映射、梯度和叶内更新。

\subsection{为什么拟合负残差方向有帮助}
对当前预测向量 $F$ 和任意修正向量 $h$，平方误差的一阶变化为
\[
J(F+\eta h)-J(F)=-\frac{\eta}{n}\sum_i (y_i-F_i)h_i+\frac{\eta^2}{2n}\sum_i h_i^2
\]
，其中 $J$ 采用半MSE。若 $h$ 与残差的内积为正，足够小的正步长可以使第一项主导而降低目标。第二项提醒我们，方向相关仍不足以允许任意大步长。这是第3章“小步沿下降方向”的思想在函数预测值上的具体形式。

一棵浅树拟合残差，就是在允许的函数集合中寻找能与残差结构一致的方向。若树直接逐点记住残差而学习率为1，可以使训练误差极低，但这正可能失去平滑和泛化能力。限制每轮学习器，并非故意降低性能，而是控制拟合过程的方式。

\section{从残差到负梯度的直观认识}

\textbf{负梯度}是损失局部下降的方向。以一条记录、当前预测$F=2$和目标$y=3$为例，损失 $L=\tfrac12(y-F)^2$ 对$F$的导数为$F-y=-1$，负梯度为$+1$。它告诉我们当前预测应该增大。另一条记录若目标为1，负梯度就是$-1$，预测应减小。树用输入特征学习哪些记录需要怎样的修正，这把逐样本的下降建议变成可用于新样本的函数。

更一般地，设当前整体模型为 $F_{m-1}$，训练损失为 $L(y_i,F(\mathbf x_i))$。在每个训练样本的当前预测处计算损失对预测的导数，并取负号：
\begin{equation}
r_{im}=-\left.\frac{\partial L(y_i,F)}{\partial F}\right|_{F=F_{m-1}(\mathbf x_i)}.
\end{equation}
新学习器尝试近似这个希望降低损失的方向。对于平方损失，负梯度与残差成比例；对于其他损失，得到的是相应的伪残差或负梯度目标。因此，“提升=永远直接拟合原始残差”是不准确的。

这里与第3章的梯度下降相通，但更新对象不完全相同。第3章直接更新一个参数向量；梯度提升从函数空间角度逐步增加学习器。理解这一点，学生就能把“梯度”作为一般优化概念，而不是只会使用某一行代码。

\subsection{概率分类的修正量不等于类别错误标记}
第6章二分类逻辑回归把得分 $F_i$ 映射为 $p_i=\sigma(F_i)$，交叉熵对得分的导数是 $p_i-y_i$。因此此处负梯度为 $y_i-p_i$。它不是 $y_i-\hat y_i$，因为 $\hat y_i$ 是阈值化后的类别，丢失了概率信息。

例如两条真实正类记录，概率分别为0.51和0.99。在0.5阈值下都分类正确，类别错误都为0；负梯度却分别为0.49和0.01，前者仍有更强的改善需求。若真实类别为0而概率为0.99，负梯度约 $-0.99$，说明应明显降低得分。这样，概率损失会区分“勉强判断正确”和“强烈支持正确类别”。

为了形成可应用于新输入的修正，需要以特征为输入拟合这些负梯度。拟合后的树叶值还可根据原损失作进一步优化。某些实现使用二阶信息：在叶中把损失近似展开为当前值加线性项和二次项，得到与梯度和曲率相关的修正。它不再只是对 $y-p$ 做算术平均。理解一般方向后，仍须依据具体软件目标解释叶值。

\subsection{一个叶的二阶修正示意}
令叶内梯度和为 $G=\sum_i g_i$，曲率和为 $H=\sum_i h_i$，在二次近似并加入 $\lambda c^2/2$ 惩罚时，要选择叶修正 $c$，近似目标为 $Gc+(H+\lambda)c^2/2$。若 $H+\lambda>0$，最小值在
\[
c^*=-\frac{G}{H+\lambda}.
\]
例如 $G=-2,H=4,\lambda=1$，修正为0.4；不加惩罚则为0.5。此例解释二阶信息和正则化怎样影响叶值，不声称所有梯度提升实现都逐字使用同一公式。对远离当前模型的大步修正，二次近似本身也可能不准确。

\subsection{与样本重加权式提升作一个区分}
AdaBoost等方法常通过增加难样本的训练权重，组织后续弱分类器。以二分类经典形式为例，若当前归一化权重下错误率为0.2，分类器权重可写成 $\tfrac12\log[(1-0.2)/0.2]=\log2$；误分样本的未归一化权重乘2，正确样本乘1/2，二者相对权重改变4倍。它突出难样本，并不等同于本节平方残差拟合。

这一小例用于认识不同提升思想的联系和区别，不增加一个大规模必修算法实验。实际使用哪一种软件类，还应核对其版本支持、目标定义和类别处理，不依据旧接口名称猜测内部算法。

\section{学习率、迭代次数与树复杂度}
较小学习率通常要求更多轮才能产生明显变化；较大学习率可能快速改善训练目标，也可能引起验证表现下降。增加 \texttt{n\_estimators} 往往使训练损失继续下降，却不保证验证误差改进。单棵树的深度和叶节点大小决定每一轮可以表达多复杂的修正，三者共同影响拟合能力。

\begin{lstlisting}[caption={住房回归：有限规模的提升参数比较}]
from sklearn.ensemble import GradientBoostingRegressor
from sklearn.model_selection import KFold, cross_val_score

cv = KFold(n_splits=3, shuffle=True, random_state=7)
for lr, trees in [(0.05, 100), (0.05, 250),
                  (0.1, 100), (0.1, 250)]:
    model = GradientBoostingRegressor(
        learning_rate=lr, n_estimators=trees,
        max_depth=2, random_state=42)
    scores = cross_val_score(model, X_dev, y_dev,
                             scoring="neg_mean_absolute_error", cv=cv)
    print(lr, trees, -scores.mean())
\end{lstlisting}

这段代码承接第5章准备的住房开发集，运行前需要先完成数据缓存。评价使用负 MAE 只是 sklearn 为统一“越大越好”接口采用的符号约定，输出时转换回正的 MAE 便于解释。若进行大量参数搜索，应减少重复计算并明确计算预算。

\subsection{学习率与轮数为何不能完全互换}
把学习率减半、轮数加倍，有时可以获得相似的总修正规模，但两者并不严格等价。每轮树是根据\emph{当时}模型重新学习的；较小步长改变后续残差及候选划分，因此会沿不同路径前进。树深也决定哪些交互能在一轮内表达，不能只用“学习率乘轮数”描述模型。

较低训练损失可能来自更多轮、更深树、较少叶样本，也可能来自错误的数据泄漏。训练曲线先用于确认程序是否按目标学习，验证曲线再用于判断这些额外拟合是否有用。即使训练目标持续下降，验证目标也可能很早达到平台，或在小样本下反复起伏。

\subsection{实际曲线与有限搜索}
本轮采用第5章共同的人工响应数据，180条训练、60条验证，固定深度2，比较学习率0.03、0.1、0.3及各自前若干轮。按验证RMSE选择的轮次如下：
\begin{center}\small
\begin{tabular}{rrrr}\toprule
学习率&选中轮数&该轮训练RMSE&该轮验证RMSE\\\midrule
0.03&218&0.2365&0.3323\\
0.10&116&0.1931&0.3134\\
0.30&60&0.1624&0.3112\\\bottomrule
\end{tabular}\end{center}
\fig{v3_ch10_stages.pdf}{人工响应数据上的逐轮验证误差。选择轮数已经使用这些验证记录，图中最低值属于开发结果，不是最终测试结果。}

较大学习率在本次设置中用较少轮达到较低验证误差，但0.1与0.3的差距很小。不能只凭这一个划分把0.3宣布为普遍最优。也不能为了证明“小学习率更好”而只展示它更平滑的一段曲线。分析应同时给出轮数、训练程度和验证范围。

\subsection{早停的含义与完整流程}
\textbf{早停}根据验证表现或指定准则，在达到最大预算之前停止增加学习器。它是一种模型选择操作，因为停止位置决定最终复杂度。若反复看最终测试误差决定停在哪里，测试已承担验证用途。使用内部验证比例时，应理解哪些记录用于建树、哪些用于早停；若数据有对象分组或时间结构，内部随机切分未必符合科学目标。

一种透明做法是在外部开发划分上训练到固定最大轮数，利用 \texttt{staged\_predict} 依次计算验证误差，按事先声明规则选择轮数，然后在全部开发数据上用该轮数重新拟合。若特征转换需要拟合，也必须在相应训练范围内学习。完成方案选择后才做一次最终测试。

\section{sklearn 梯度提升实践}
对于分类可以使用 \texttt{GradientBoostingClassifier}，对回归使用 \texttt{GradientBoostingRegressor}。更大数据量也可考虑直方图梯度提升实现 \texttt{HistGradientBoostingRegressor} 等。选择算法前先检查数据规模、特征类型、缺失数据和资源限制。不同实现可能有不同的树构造和正则化控制方式，不能假设同名参数在各包中完全等价。

课堂主实验固定住房数据和回归评价方案，与常数基线、线性回归、单树和随机森林比较。先选择一种简单的深度配置，再比较少量学习率和迭代次数。即便提升方法最终分数更高，也应检查它在高目标值、稀少工况和不同地区的误差是否改善。

\subsection{把逐轮诊断落实为可检查的代码}
以下片段承接同一训练与验证表，先记录每轮验证误差，再确定候选轮次。它不会在每轮重新训练一个独立模型，而是查看同一个加法模型的前缀。
\begin{lstlisting}[caption={在开发验证集上选择轮数，并明确这是调参}]
import numpy as np
from sklearn.ensemble import GradientBoostingRegressor
from sklearn.metrics import root_mean_squared_error
model = GradientBoostingRegressor(
    n_estimators=250, learning_rate=.1, max_depth=2,
    random_state=42)
model.fit(X_tr, y_tr)
val_error = np.array([
    root_mean_squared_error(y_val, prediction)
    for prediction in model.staged_predict(X_val)
])
best_round = int(np.argmin(val_error))+1
print(best_round, val_error[best_round-1])
\end{lstlisting}

验证噪声较大时，可以选择与最低误差接近但轮数更少的模型，前提是接近的判据在查看最终测试之前固定。观察到一次局部小下降，不必立即把它解释为新的物理结构。模型需要的轮数还随样本量、特征表达和目标噪声改变。

\subsection{直方图实现改变的是搜索近似与工程代价}
直方图梯度提升先把连续特征映射到有限箱，再在箱边界上比较候选划分。与逐个原始数值中点搜索相比，候选数量减少，能够带来不同的计算与存储代价；代价是输入值在搜索时被离散化。箱边界、缺失处理和正则化设置会影响训练，不能把所有提升实现当作完全相同模型的纯加速版本。

本书先用普通梯度提升解释机制，再把直方图版本作为较大表格数据的实现比较。比较时固定任务、数据划分和可用字段，同时记录各自参数含义、实际运行时间及误差。缺失值支持或类别特征支持属于具体版本接口，需要以随书环境与官方文档为准，不根据名称推断。

\section{与随机森林的比较}
两种方法都以树为基础，但组织结构不同。随机森林通过多次重采样与特征随机性形成有差异的树，再综合预测；梯度提升有顺序地增加修正模型。随机森林常对参数选择相对稳健，提升可能对学习率与迭代轮数更加敏感，但具体表现取决于数据与实现。



模型选择不必寻找唯一“最好”的方法：对一些科学任务，训练成本、解释便利性、推理延迟和可复现程度同样重要。应将这些因素与评价指标一起记录。

\subsection{相同树组件，为什么训练结构不同}
森林可以在同一批原始目标上分别训练多棵树，各树通过抽样和特征候选形成差异；提升则把前一轮预测带入下一轮目标。森林预测是并列聚合，提升预测是初值加修正。森林中新增一棵树主要改变平均，提升中新增一棵树则针对当前模型的不足而选择方向。

这种区别影响并行与诊断。森林的树之间较容易并行训练；提升的轮次通常依赖前轮，但一轮内部仍可并行处理特征或数据。不能由此得出“森林总比提升快”：具体树深、数据量、实现与硬件都需要实测。统计上也不存在“森林只降方差、提升只降偏差”的绝对分工，两者都可能改变多个误差来源。

\subsection{共同人工数据上的一次完整比较}
本轮固定第5章人工响应数据与180/60训练验证划分，采用一棵限深树、一个随机森林、一个固定配置提升模型，结果如下。该表的提升模型不是上一节为每个学习率分别挑出的最低点，配置在实验脚本中单独记录，避免混淆比较口径。
\begin{center}\small
\begin{tabular}{lrr}\toprule
模型&训练RMSE&验证RMSE\\\midrule
单棵树&0.3128&0.4304\\
随机森林&0.2279&0.3413\\
梯度提升&0.1900&0.3140\\\bottomrule
\end{tabular}\end{center}
在这一人工任务上，两个集成均优于单树，提升验证误差最低。结果支持的是本数据、本划分和这些配置之间的比较，不支持跨任务统一排名。分析下一步应检查残差是否集中在输入边缘、是否漏掉某个振荡结构，以及模型在保留范围外如何表现，而不是继续无休止地搜索更多树数。

与线性回归相比，树集成的灵活性可能更容易拟合非线性，但也失去简单系数的直接解释。对近似线性、噪声较大的小样本，线性基线可能已经足够。是否采用复杂模型要结合误差改善、验证稳定性和应用成本共同决定。

\section{XGBoost 扩展与专业任务}

\textbf{XGBoost}是一个独立的软件项目名称，不是本章中新增加的一种任务类型。它在梯度提升思想上提供具体的目标、正则化和高效实现。比较两个软件实现时，要先确认训练目标、树的复杂度及验证方案是否相当，再解释分数差异。\textsuperscript{[12]}

XGBoost 提供高效的梯度提升树实现及与 sklearn 风格兼容的估计器接口，但不是 sklearn 核心包的一部分。本书先使用 sklearn 建立清晰的提升概念，再在附录中提供 XGBoost 的可选安装与基本对照。不要为了追求某个软件名称而跳过模型评价。

专业方向可使用 Airfoil 继续回归比较，检验不同频率与流速条件下的误差；提高项目可使用超导临界温度数据。

\label{sec:superconduct-intro}
\textbf{超导临界温度}指材料在指定条件下进入超导状态的转变温度，本项目以该温度作为数值预测目标。所用\textbf{超导材料数据集（Superconductivty Data，保留数据源的拼写）}包含21\,263条记录，主表由81个材料描述特征和临界温度组成，另有文件提供化学组成信息。它们不是81次温度测量，而是由元素及组成信息整理出的描述量。数据字段与来源应依据原始说明和论文核对。\textsuperscript{[18]}

一条记录的样本单位是材料记录；化学式相同的多条记录可能并不独立。读取辅助组成文件时还须注意，其中也包含目标温度，不能把它在合并时重复加入输入。该项目用于训练材料建模和分组评价，不能把预测误差很小直接等同于找到了新的超导机制。
后者包含不同化学组成与材料特征，样本分组和可能的材料族外推比单一高分更值得讨论。若准备数据时发现多条记录共享配方或物理组成，需明确怎样进行去重或分组验证。

\subsection{软件扩展前，先建立可以复用的比较协议}
引入XGBoost时，保留同一输入表、对象标识、开发划分、主要指标和错误分析，而不是重新随机拆分一次再与旧分数比较。明确所用目标是连续回归还是类别概率，叶数或深度怎样约束，训练过程中是否使用了独立早停数据。相同名称的学习率或正则化参数也可能具有不同的归一化约定，不能只抄数值。

附录中的XGBoost实验是可选依赖，不是运行全书的前提。核心梯度提升实验使用已锁定环境中的sklearn完成。安装失败或没有该包时，程序应清楚报告可选项未执行，而不能把另一个模型的结果标成XGBoost。

\subsection{从较低残差走向科学解释}
风洞声学数据中，增加树数以后整体RMSE下降，不代表所有频率与流速区间都改善。应按可解释工况分组，检查模型是否只在占比大的常见条件上更好。对于某些未覆盖工况，模型仍只是把输入送进已有叶，再叠加有限修正；复杂度不能创造缺失的实验信息。

超导材料数据则需要先判断“新材料”的含义：是同一组成的重复记录、已知材料族的新配比，还是完全不同的组成族？三者对应不同划分与泛化难度。一个在随机行划分中表现很好的模型，可能只是利用与训练组成高度相近的记录；要研究组成族外推，必须以适当分组建立评价。程序从输入描述量预测临界温度，并没有直接求解超导微观机制。

本章最终应留下两种能力：能在小数据上逐轮核对修正，解释为什么损失下降；能在实际任务中建立有限、透明的调参和误差分析。前者帮助识别算法误解，后者避免把训练改善直接写成科学发现。

\chapterpractice{先在模拟一维数据上可视化前几轮提升，再在住房开发集上比较 GradientBoostingRegressor 与 RandomForestRegressor，用小网格研究学习率与树数。专业迁移在 Airfoil 或经核对的超导数据上开展。}
\chapterexercise
\begin{enumerate}
\item 为什么平方损失下拟合残差能够解释梯度提升？换成交叉熵时是否仍然直接拟合 $y-\hat y$？
\item 同时提高学习率和树数，能否仅凭更低训练损失证明泛化改善？
\item 随机森林与提升树都由很多树组成，为什么不能把二者理解为同一种训练策略？
\item 在超导材料预测中，为什么同配方重复记录可能影响随机划分的结果？
\end{enumerate}

\section*{深化练习与方法检验}
本组题强调解释、计算与验证；D编号在本章内独立计数。计算题应写出所用定义和中间步骤，实验题保留实际输出，包括不支持预期的结果。
\begin{enumerate}[label=D\arabic*.]
\item 对目标 $(0,1,2,4)$，从均值1.75开始，以学习率0.5加入正文的第一棵残差树。写出新预测与MSE。说明为什么“基学习器预测残差”并不意味着最终输出就是残差。
\item 对单个二分类样本，当前得分0、类别1，求交叉熵对得分的负梯度。若误把类别数字减去得分作为修正目标，会得到什么，为什么不同？
\item 两个提升配置，一个用学习率0.05训练200棵树，另一个用0.1训练100棵树。乘积相同是否意味着最终模型相同？说明理由并给出验证设计。
\item 在同一验证集上从300轮中选择最优停止轮数后，这个最低验证RMSE能否当作最终独立测试结果？说明正确的后续步骤。
\end{enumerate}
